\documentclass[11pt,a4paper]{article}

\usepackage{booktabs}
\usepackage[utf8]{inputenc}
\sloppy
\usepackage[T1]{fontenc}

% allow line breaks in long URLs
\PassOptionsToPackage{hyphens}{url}

\usepackage{geometry}
\geometry{margin=2.5cm}

\usepackage{hyperref}
\hypersetup{
    breaklinks=true,
    colorlinks=true,
    urlcolor=blue
}

% -------------------------
% ---- STORAGE LINKS -------
% -------------------------

% Local master data folder
\newcommand{\localdata}{C:/Users/aliad/Documents/osteomin}

% UiO server (primary lab storage)
\newcommand{\uioserverdata}{/mn/fys-server1/kondfys/aliaslan/data/}

% Cloud copies (Google Drive + OneDrive)
\newcommand{\googledata}{https://drive.google.com/drive/folders/1QbGXquTdg0ZWvXDNWVM-L-pLeBURobRO?usp=sharing}
\newcommand{\onedrivedata}{https://uio-my.sharepoint.com/:f:/g/personal/aliadem_uio_no/IgBenICVupn0RIL3Lm9zlAd3ASNnG41zc4TH_MuMsQkTJFc}

% -------------------------
% ---- CODE LINKS ----------
% -------------------------

\newcommand{\gitmain}{https://github.com/aliaslandemir/osteomin}
\newcommand{\gitanalysis}{https://github.com/aliaslandemir/osteomin/tree/main/endpointanalysis}

% -------------------------
% ---- MANUSCRIPT LINKS ----
% -------------------------

% Overleaf projects
\newcommand{\ovlmain}{https://www.overleaf.com/project/68d16219d9b6ab7685e7c2b5}
\newcommand{\ovlsupp}{https://www.overleaf.com/project/68d16219d9b6ab7685e7c2b5}

% -------------------------



\title{OsteoMin Project: Data, Code, and Manuscript Map}
\author{Ali Aslan Demir}
\date{\today}

\begin{document}

\maketitle

\section{Purpose}
Paths for data, analysis codes, manuscripts, and more used in the OsteoMin project. This document contains where are the data, codes, manuscript. How to run endpointanalysis codes. And also; name method for microscopy images.

\section{Experimental Data}

\subsection{Local Master Folder}
\begin{itemize}
  \item Local primary data: \texttt{\localdata}
  \item Synchronized to both Google Drive (Computers/My Laptop) and UiO OneDrive.
\end{itemize}

\subsection{UiO Server (Primary Laboratory Copy)}
\begin{itemize}
  \item Server path: \texttt{\uioserverdata}
  \item Contains ALP, confocal (Collagen~I), brightfield time-series, ARS, and metadata folders.
\end{itemize}

\subsection{Cloud Copies}
\begin{itemize}
  \item Google Drive (Computers sync): \\[2pt]\url{\googledata}
  \item UiO OneDrive (this one is problematic to synch! I will upload manually to cloud to be sure. Use Google Drive and Github primarily.): \\[2pt]\url{\onedrivedata}
\end{itemize}

\section{Code and Simulations}

\subsection{GitHub Repositories}
\begin{itemize}
  \item OsteoMin Simulation Model: \\[2pt]\url{\gitmain}
  \item Endpoint analysis scripts: \\[2pt]\url{\gitanalysis}
\end{itemize}


\section{Manuscripts}

\subsection{Overleaf Projects}
\begin{itemize}
  \item Main manuscript: \\[2pt]\url{\ovlmain}
  \item Supplementary material
  \item osteomin guide
  \item fitapproach
  \item bibliograpghy
\end{itemize}





\newpage

\section{How to Run Codes}

% -------------------------------------------------------

\subsection{Actin Analysis}

Use actin channel of confocal via \texttt{actin\_analysis.py}  
(see also \texttt{actin\_an\_meta.md}).  
Input folder: \texttt{osteomin/data/confocal}.

\subsection*{Command}

\begin{center}
\begin{minipage}{0.9\linewidth}
\begin{verbatim}
python actin_analysis.py \
    --input ../data/input_images \
    --out actin_results
\end{verbatim}
\end{minipage}
\end{center}

% -------------------------------------------------------

\subsection{Collagen Analysis}

Use collagen IF data with \texttt{collagen\_analysis.py}  
(see also \texttt{collagen\_an\_meta.md}).  
Point \texttt{--input} at your confocal folder and choose any \texttt{--out} root.

\subsection*{Command}

\begin{center}
\begin{minipage}{0.9\linewidth}
\begin{verbatim}
python collagen_analysis.py \
    --input ../data/input_images \
    --out collagen_out
\end{verbatim}
\end{minipage}
\end{center}

% -------------------------------------------------------

\subsection{DAPI Analysis}

Use DAPI channel images via \texttt{dapi\_analysis.py}  
(see also \texttt{dapi\_an\_meta.md}).  
Provide the same \texttt{--input} folder and direct results with \texttt{--out}.

\subsection*{Command}

\begin{center}
\begin{minipage}{0.9\linewidth}
\begin{verbatim}
python dapi_analysis.py \
    --input ../data/input_images \
    --out dapi_out
\end{verbatim}
\end{minipage}
\end{center}

% -------------------------------------------------------

\subsection{Overlay Generation}

Use \texttt{overlay\_analysis.py} to create merged overlays  
(see also \texttt{overlay\_an\_meta.md}).  
Point \texttt{--input} at your confocal folder and select any \texttt{--out}.

\subsection*{Command}

\begin{center}
\begin{minipage}{0.9\linewidth}
\begin{verbatim}
python overlay_analysis.py \
    --input ../data/input_images \
    --out overlay_output \
    --pix 0.21 \
    --scalebar_um 100
\end{verbatim}
\end{minipage}
\end{center}

% -------------------------------------------------------

\subsection{ARS Analysis}

Use \texttt{data/ars} and \texttt{ars\_analysis.py}  
(see also \texttt{ars\_an\_meta.md}).  
Pass \texttt{--input} and \texttt{--out} (or \texttt{--output}) just like the confocal pipelines.

\subsection*{Command}

\begin{center}
\begin{minipage}{0.9\linewidth}
\begin{verbatim}
python ars_analysis.py \
    --input ../data/ars \
    --out ars_out \
    --lcap 65
\end{verbatim}
\end{minipage}
\end{center}


% -------------------------------------------------------

\subsection{ALP Analysis}

Use \texttt{data/alp} and \texttt{alp\_analysis.py}  
(see also \texttt{alp\_an\_meta.md}).  
Point \texttt{--input} to the Excel sheet and \texttt{--out} to the results root.

\subsection*{Command}

\begin{center}
\begin{minipage}{0.9\linewidth}
\begin{verbatim}
python alp_analysis.py \
    --input data/alp/alp_raw_data.xlsx \
    --out alp_results
\end{verbatim}
\end{minipage}
\end{center}



\subsection{OsteoMin Cellular Automaton}

The cellular automaton is run from the main OsteoMin repository using
\texttt{main.py}.  
Activate the virtual environment first:

\begin{center}
\begin{minipage}{0.9\linewidth}
\begin{verbatim}
source .venv/bin/activate
\end{verbatim}
\end{minipage}
\end{center}

\subsection*{Default Run Command}

\begin{center}
\begin{minipage}{0.95\linewidth}
\begin{verbatim}
python main.py \
    --days 30 \
    --cells 72 \
    --w 33 --h 33 \
    --boundary-mode periodic \
    --seed $(shuf -i 1-999999 -n 1) \
    --media-interval 2 \
    --aa 1.0 --bgp 1.0 \
    --random-move --move-prob 1.0 --max-moves 1 \
    --proliferation --ob-div-prob 0.0 \
    --apoptosis --apoptosis-prob 0.01 --apoptosis-osteo 0.001 \
    --osteogenic-commitment --commit-growth 0.01 --commit-delay 3 \
    --secretion --nucleation --enable-nodules \
    --spatial-every 1 \
    --out out \
    --nodule-window 4 \
    --k-ha 0.06 --k-nuc 0.025
\end{verbatim}
\end{minipage}
\end{center}

\newpage

\subsection{File naming convention and condition labels}

Microscopy files follow:
\[
\texttt{C\{n\}\_D\{day\}\_\{objective\}\_\{modality\}\_\{idx\}.tif}
\]

Example:
\texttt{C4\_D28\_4x\_BF\_04.tif}

\noindent\textbf{Components:}
\begin{itemize}
    \item \texttt{C\{n\}} — condition (C1--C8)
    \item \texttt{D\{day\}} — day (D09, D14, D28 ...)
    \item \texttt{objective} — 4x, 10x, 20x
    \item \texttt{modality} — BF, DAPI, Actin, Collagen
    \item \texttt{idx} — image/fov number
\end{itemize}

\subsubsection*{Condition $\rightarrow$ Label Mapping}

\begin{table}[htbp]
\centering
\renewcommand{\arraystretch}{1.1}
\setlength{\tabcolsep}{8pt}
\begin{tabular}{@{}l l l@{}}
\toprule
\textbf{Code} & \textbf{Condition} & \textbf{Figure label} \\
\midrule
C1 & Basal & Basal \\
C2 & Osteo & A+$\beta$GP \\
C3 & Osteo + Dex/MK4 (L/L) & A+$\beta$GP, $D_LK_L$ \\
C4 & Osteo + Dex/MK4 (L/H) & A+$\beta$GP, $D_LK_H$ \\
C5 & Osteo + Dex/MK4 (H/L) & A+$\beta$GP, $D_HK_L$ \\
C6 & Osteo + Dex/MK4 (H/H) & A+$\beta$GP, $D_HK_H$ \\
C7 & Osteo + Dex (H) & A+$\beta$GP, $D_HK_0$ \\
C8 & Osteo + MK-4 (H) & A+$\beta$GP, $D_0K_H$ \\
\bottomrule
\end{tabular}
\end{table}

\noindent
L/H indicate low or high doses; $D_0$ or $K_0$ indicate absence of Dex or MK-4.



\end{document}
